\begin{table}[H]
  \centering
  \caption{四问所用的物性与经验式汇总}\label{tab:params}
  \footnotesize
  \begin{tabularx}{\textwidth}{@{}l >{\raggedright\arraybackslash}X >{\raggedright\arraybackslash}X >{\raggedright\arraybackslash}X@{}}
    \toprule
    参数 & 第 (1) 问（附录 2） & 第 (2)(3) 问（附录 3） & 第 (4) 问（附录 4） \\
    \midrule
    密度 / (kg/m$^3$) & $820$ & $650+128C$ & $760+90C$ \\
    比热容 / (J/(kg$\cdot$K)) & $2600$ & $1450+2736C/(C+1)$ & $1850+2150C/(C+1)$ \\
    导热系数 / (W/(m$\cdot$K)) & $0.36$ & $0.21+0.38C/(C+1)$ & $0.12+0.20C/(C+1)$ \\
    扩散系数 / (m$^2$/s) & $7\times10^{-9}e^{-0.89/C}$ & $2.4\times10^{-3}e^{-0.45/C}e^{-3850/T}$ & $4.2\times10^{-4}e^{-0.30/C}e^{-3850/T}$ \\
    \midrule
    对流传热系数 $h$ / (W/(m$^2\cdot$K)) & \multicolumn{3}{c}{$25$} \\
    对流传质系数 $h_m$ / (m/s) & \multicolumn{3}{c}{$8\times10^{-7}$} \\
    初始温度 / $^\circ$C & \multicolumn{3}{c}{$28$} \\
    初始干基含水率 / (kg/kg) & \multicolumn{3}{c}{$2.55$} \\
    烘干判据 / (kg/kg) & \multicolumn{3}{c}{全域最大 $\le 0.15$} \\
    \bottomrule
  \end{tabularx}
  \par\smallskip\footnotesize 表中 $T$ 在经验式内为热力学温度，单位为 K；其余各处温度均为摄氏度。
\end{table}
